\section*{Chapter \ref{ch:m2:the-fx-options-market} --- The FX Options Market}

\begin{solution}{exo:m2:the-fx-options-market:1}
$7 + 0.2 - 0.25 = 6.95\%$ for the 25-delta euro call and $7 + 0.2 + 0.25 = 7.45\%$ for
the put: the euro put is dearer, by the half-point of the \omterm{def:m2:the-fx-options-market:rr}{risk reversal}.
\end{solution}

\begin{solution}{exo:m2:the-fx-options-market:2}
EURUSD regular deltas (premium in dollars, the quote currency); USDJPY
premium-adjusted (premium in dollars, the base currency); USDBRL forward and
premium-adjusted (an emerging-market pair, premium in dollars). The premium paid
in the base currency is a position in it, which changes the hedge and so the
delta that defines each quoted strike.
\end{solution}

\begin{solution}{exo:m2:the-fx-options-market:3}
$600\,000 - 73\,669 = 526\,331$ euros: a \omterm{def:m2:the-fx-options-market:delta}{premium-adjusted delta} of 52.63\%.
\end{solution}

\begin{solution}{exo:m2:the-fx-options-market:4}
Regular \omterm{def:m2:the-fx-options-market:delta}{spot deltas}: call 1.17904, put 1.12357; \omterm{def:m2:the-fx-options-market:delta}{forward deltas}: 1.17920 and
1.12341.
\end{solution}

\begin{solution}{exo:m2:the-fx-options-market:5}
With regular deltas the call delta $N(d_1)$ equals the put's $N(-d_1)$ when
$d_1 = 0$, that is $K = Fe^{\sigma^2T/2}$, above the forward. The \omterm{def:m2:the-fx-options-market:delta}{premium-adjusted
delta} uses $N(d_2)$ scaled by $K/F$, and neutrality requires $d_2 = 0$, that is
$K = Fe^{-\sigma^2T/2}$, below it.
\end{solution}

\begin{solution}{exo:m2:the-fx-options-market:6}
It pays nothing in the paths where spot touches the barrier before expiry, some
of which would have ended in the money. The difference is large when the barrier
is close to spot, volatility is high and the expiry long, since the touch is
then likely.
\end{solution}

\begin{solution}{exo:m2:the-fx-options-market:7}
With \omterm{def:m2:the-fx-options-market:delta}{forward deltas} the call strike is 160.90 and the put 150.65, against 160.85 and
150.71 with regular \omterm{def:m2:the-fx-options-market:delta}{spot deltas} and 160.68 and 150.52 with premium-adjusted \omterm{def:m2:the-fx-options-market:delta}{spot
deltas}.
\end{solution}

\begin{solution}{exo:m2:the-fx-options-market:8}
A \omterm{def:m2:the-fx-options-market:rr}{risk reversal} is a price, not a forecast: dollar puts are dearer than calls
because more buyers want protection against a fall, and dealers who sell it
charge for the risk of a sharp move that way. It says the market pays more to
insure against a dollar fall; the forward, not the smile, carries the market's
expected rate under the pricing measure.
\end{solution}

\begin{solution}{pb:m2:the-fx-options-market:1}
\textbf{1.} 183.8 pips of notional (0.01838 dollars per euro).
\textbf{2.} 165.2 pips; the company saves 18.6 pips, about USD~930\,000 on EUR~500
million.
\textbf{3.} 0.657; the dealer holds about EUR~328.7 million.
\textbf{4.} Near the barrier the knock-out's value falls to zero much faster than
the vanilla's as spot falls, so its value is more sensitive to spot.
\textbf{5.} 0.631 just above 1.1200, against 0.298 for the vanilla.
\textbf{6.} About EUR~315.5 million, long euros.
\textbf{7.} Sell all of it, EUR~315.5 million, at the touch.
\textbf{8.} It adds a large sale at the moment spot is falling through the barrier,
pushing it lower.
\textbf{9.} Buy euros near 1.1200, to keep the barrier from trading and the option
alive.
\textbf{10.} If they know a large hedge unwinds there, pushing spot to it may
trigger the sale and a move they can profit from.
\textbf{11.} By selling part of the hedge in advance as spot approaches, or by
buying vanilla options that offset the barrier's delta jump.
\textbf{12.} It gives up hedge accuracy: if spot turns back, the dealer is short
the delta it sold early; the options cost premium.
\textbf{13.} Gamma near a knock-out barrier is large and changes sign across it:
small moves require large rebalancing trades, each paying a spread.
\textbf{14.} With a smile the volatility at the barrier (a euro put's) is higher,
the probability of touching higher, and the knock-out cheaper than with the flat
volatility.
\textbf{15.} That the protection disappears exactly when spot falls through 1.1200,
and that its discount is the price of that risk.
\textbf{16.} Hedging an exposure is legitimate; trading to prevent or cause a
touch in order to benefit from it, using client information, is not, and falls
under the rules against manipulation and the \omterm{def:m2:fx-spot-microstructure:code}{FX Global Code}.
\textbf{17.} Clients choose round levels, which are easy to remember and watch.
\textbf{18.} Clients' barrier levels and sizes, which reveal where others could push
the market.
\textbf{19.} Named result: \emph{the barrier defence} sets the dealer's sale at the
touch at about EUR~315.5 million, a delta of 0.631 on EUR~500 million, all at
1.1200.
\textbf{20.} Because the hedges of options that vanish or appear at a level must be
unwound or put on there, in size and at once.
\end{solution}

\begin{solution}{iq:m2:the-fx-options-market:1}
The ATM volatility is the level; the 25-delta \omterm{def:m2:the-fx-options-market:rr}{risk reversal}, the call's volatility
minus the put's, is the skew: which side the market pays more to insure; the
\omterm{def:m2:the-fx-options-market:rr}{butterfly}, the wings' average over ATM, is the curvature: how fat the tails are
priced. Together they give three points of the smile.
\iqlookfor{level, skew and curvature.}
\end{solution}

\begin{solution}{iq:m2:the-fx-options-market:2}
The delta corrected for a premium paid in the base currency, which is itself a
position in that currency: $\phi(K/F)N(\phi d_2)$ in forward terms, or discounted
for spot. It is the convention for pairs whose premium is paid in the base
currency, such as USDJPY or EURJPY.
\iqlookfor{why the premium changes the hedge, and which pairs.}
\end{solution}

\begin{solution}{iq:m2:the-fx-options-market:3}
Convert the quotes to the put's volatility, by the simplified formula or a
calibration to the market \omterm{def:m2:the-fx-options-market:rr}{strangle}; then solve for the strike at which a put with
that volatility has a delta of $-0.25$ in the pair's convention: closed form for
regular deltas, a monotone root search for premium-adjusted ones.
\iqlookfor{the order: volatility first, then strike, in the right delta.}
\end{solution}

\begin{solution}{iq:m2:the-fx-options-market:4}
Estimate the delta that will be released at the touch and the gamma on the way;
sell part of the hedge early or buy options that offset the jump, within limits;
keep the barrier confidential; do not trade to cause or prevent the touch; if
the touch comes, execute the unwind as planned rather than chase.
\iqlookfor{preparation, not manipulation.}
\end{solution}

\begin{solution}{iq:m2:the-fx-options-market:5}
Because the straddle at that strike has no delta, so it trades as a pure
volatility instrument and needs no hedge on the trade; the forward strike would
leave a small delta. The convention reflects what traders use at-the-money
options for: trading volatility.
\iqlookfor{delta neutrality as the purpose.}
\end{solution}

\begin{solution}{iq:m2:the-fx-options-market:6}
Store per pair its conventions (delta type by tenor, ATM definition, premium
currency); ingest ATM, RR and BF by tenor with timestamps; calibrate a smile per
expiry exactly to the quotes; interpolate in time on total variance; serve
volatilities by strike or delta with the convention stated; version every surface;
check for arbitrage (calendar and \omterm{def:m2:the-fx-options-market:rr}{butterfly}) and stale quotes; test against
hand-built cases.
\iqlookfor{conventions as data, exact calibration, arbitrage checks.}
\end{solution}
